\documentclass[10pt]{article}
% ---------------------------------------------------------------------------
%  ccint project wrap-up — preamble
%  Compiles with pdfLaTeX (Overleaf default; no configuration needed).
% ---------------------------------------------------------------------------
\usepackage[T1]{fontenc}
\usepackage[utf8]{inputenc}
\usepackage{lmodern}
\usepackage[letterpaper,margin=1in,headheight=14pt]{geometry}
\usepackage{microtype}
\setlength{\emergencystretch}{3em}   % 允许必要时额外伸缩，消除长代码标识符造成的溢出

\usepackage{booktabs}
\usepackage{tabularx}
\usepackage{longtable}
\usepackage{array}
\usepackage{multirow}
\usepackage{amssymb}
\usepackage{amsmath}
\usepackage{enumitem}
\usepackage{fancyhdr}
\usepackage{titlesec}
\usepackage{graphicx}
\graphicspath{{figures/}}
\usepackage{tikz}
\usetikzlibrary{arrows.meta,positioning,calc,fit,backgrounds}
\usepackage[skins,breakable]{tcolorbox}

\usepackage[dvipsnames]{xcolor}
\definecolor{ink}{HTML}{1A1A1A}
\definecolor{accent}{HTML}{1F4E79}
\definecolor{warn}{HTML}{9C4221}
\definecolor{good}{HTML}{22543D}
\definecolor{rule}{HTML}{C8CDD4}
\definecolor{softbg}{HTML}{F4F6F8}
\definecolor{accentbg}{HTML}{EAF0F6}

\usepackage[font=small,labelfont={bf,color=accent},skip=8pt]{caption}

\usepackage[colorlinks=true,linkcolor=accent,urlcolor=accent,citecolor=accent]{hyperref}

% --- headers ---------------------------------------------------------------
\pagestyle{fancy}
\fancyhf{}
\renewcommand{\headrulewidth}{0.4pt}
\fancyhead[L]{\footnotesize\textcolor{accent}{ccint — Canada Cyber Social Intelligence}}
\fancyhead[R]{\footnotesize\thepage}

% --- section styling -------------------------------------------------------
\titleformat{\section}{\Large\bfseries\color{accent}}{\thesection}{0.7em}{}
\titleformat{\subsection}{\large\bfseries\color{ink}}{\thesubsection}{0.6em}{}
\titleformat{\subsubsection}{\normalsize\bfseries\itshape\color{ink}}{}{0em}{}
\titlespacing*{\section}{0pt}{2.2ex plus 1ex minus .2ex}{1.2ex plus .2ex}

% --- table helpers ---------------------------------------------------------
\renewcommand{\arraystretch}{1.18}
\newcolumntype{R}{>{\raggedleft\arraybackslash}X}
\newcolumntype{L}{>{\raggedright\arraybackslash}X}
\newcommand{\hd}[1]{\textbf{#1}}

% --- callout boxes ---------------------------------------------------------
\newtcolorbox{keybox}[1][]{
  enhanced, breakable, colback=accentbg, colframe=accent,
  boxrule=0pt, leftrule=3pt, arc=1pt, left=8pt, right=8pt, top=6pt, bottom=6pt,
  #1}
\newtcolorbox{warnbox}[1][]{
  enhanced, breakable, colback=softbg, colframe=warn,
  boxrule=0pt, leftrule=3pt, arc=1pt, left=8pt, right=8pt, top=6pt, bottom=6pt,
  #1}
\newtcolorbox{plainbox}[1][]{
  enhanced, breakable, colback=softbg, colframe=rule,
  boxrule=0.4pt, arc=1pt, left=8pt, right=8pt, top=6pt, bottom=6pt,
  #1}

% --- inline semantics ------------------------------------------------------
\newcommand{\code}[1]{\texttt{\small #1}}
\newcommand{\yes}{\textcolor{good}{\checkmark}}
\newcommand{\no}{\textcolor{warn}{$\times$}}
\newcommand{\noise}{\textcolor{warn}{noise}}

% --- verbatim-ish code block ----------------------------------------------
\usepackage{fancyvrb}
\DefineVerbatimEnvironment{shell}{Verbatim}
  {fontsize=\small,frame=single,rulecolor=\color{rule},framesep=6pt,
   formatcom=\color{ink}}


% --- 两页版：收紧版面 ---------------------------------------------------
\geometry{margin=0.82in}
\setlength{\parskip}{0.28em}
\setlength{\parindent}{0pt}
\renewcommand{\arraystretch}{1.10}
\titlespacing*{\section}{0pt}{1.4ex plus .3ex minus .2ex}{0.7ex}
\titleformat{\section}{\large\bfseries\color{accent}}{\thesection.}{0.5em}{}
\pagestyle{fancy}
\fancyhf{}
\renewcommand{\headrulewidth}{0pt}
\fancyfoot[C]{\footnotesize\thepage\ / 2}

\begin{document}

\begin{center}
{\LARGE\bfseries ccint --- Canada Cyber Social Intelligence}\\[0.3em]
{\normalsize Two-page summary \quad$\cdot$\quad full write-up available separately}\\[0.4em]
\textcolor{accent}{\rule{\textwidth}{0.9pt}}\\[0.2em]
{\footnotesize Status 2026-09-21 \quad$\cdot$\quad Corpus 2026-08-22 $\to$ 2026-09-21 (30 days)}
\end{center}
\vspace{-0.3em}

% ==========================================================================
\section{What it is}

A longitudinal system that continuously collects Canada-related cybersecurity
discussion from Bluesky, stores it as a timestamped corpus, tests it for
temporal change, and has an LLM agent produce evidence-grounded analysis of
what survives that test.

\vspace{0.2em}
\begin{minipage}[t]{0.47\textwidth}
\small
\begin{tabular}{@{}lr@{}}
\toprule
Posts collected & 90{,}503 \\
Distinct authors & 26{,}810 \\
Canada $\times$ cyber relevant & 650 (0.72\%) \\
\textbf{Relevant posts / day} & \textbf{21.3} \\
Collection runs & 78 \\
Tests & 191 passing \\
\bottomrule
\end{tabular}
\end{minipage}\hfill
\begin{minipage}[t]{0.50\textwidth}
\vspace{-1.2em}
\begin{keybox}
\small
\textbf{The headline result is a negative one.} No topic in the current window
is distinguishable from random noise --- including the one that looks like
$+56\%$ growth.
\end{keybox}
\end{minipage}

% ==========================================================================
\section{Why Bluesky, and what that costs}

The corpus accrues with wall-clock time: a day not collected is a day
permanently lost. So the question was not \emph{which platform is best} but
\emph{which one can I start collecting from today}. X/Twitter is gated behind
paid research access; Mastodon has no global search, so corpus completeness is
undefinable; Reddit is viable and remains the best second source. Bluesky offers
open API access with no approval step and verified history at least three months
back.

\textbf{It was chosen for API availability, not representativeness.} Its user
base is skewed and much smaller than where most Canadians actually talk. Nothing here should be read as
``what Canadians think'' --- it is ``what appears on Bluesky,'' and that gap is
real and unmeasured.

% ==========================================================================
\section{The design decision that determines everything}

\begin{keybox}
\textbf{The collector searches 34 cybersecurity terms. Zero of them mention
Canada.} Everything is stored; Canada relevance is decided \emph{afterwards},
offline, against the stored corpus.
\end{keybox}

The asymmetry is deliberate. A \emph{collection} filter is irreversible --- posts
never fetched cannot be recovered. A \emph{relevance} filter is reversible, as
long as the raw payload survives. Cyber vocabulary is stable enough to bound
collection; ``what counts as Canada-related'' is exactly the definition most
likely to change, so it belongs on the recomputable side.

\textbf{This paid off three times in one day.} \code{GTA} turned out to mean
\emph{Grand Theft Auto} (18.1\% of ``relevant'' was video-game news);
\code{CRA} meant the EU \emph{Cyber Resilience Act}; \code{CSIS} meant a US
think tank. All three were fixed by writing a new label version --- \textbf{no
re-collection}. Had ``canada'' been a search term, those errors would have been
frozen into the corpus permanently. The cost is 139$\times$ storage overhead,
which is the cheap side of that trade.

% ==========================================================================
\section{What the data actually says}

Window comparison answers ``which topic changed most,'' not ``is that change
larger than noise.'' A permutation test closes the gap: each author's entire
posting history is circularly shifted across the 14-day span, preserving post
counts, topic mix and burst structure while destroying calendar alignment.

\vspace{0.2em}
\begin{center}
\small
\begin{tabular}{@{}lrrrl@{}}
\toprule
\hd{topic} & \hd{$\Delta$} & \hd{growth} & \hd{$p$ (FWER)} & \hd{verdict} \\
\midrule
\code{data\_breach}         & $+10$ & $+56\%$  & 0.718 & \noise \\
\code{other}                & $+4$  & $+5\%$   & 0.991 & \noise \\
\code{malware\_infostealer} & $+3$  & $+100\%$ & 0.998 & \noise \\
\code{ransomware}           & $-5$  & $-20\%$  & 0.969 & \noise \\
\bottomrule
\end{tabular}
\end{center}

\begin{keybox}
Under the null, random fluctuation alone produces a swing of $\pm28$ in
\emph{some} topic. The largest observed change is $+10$. $p$-values are
corrected for having picked the largest of 10 topics to narrate.
\end{keybox}

\textbf{$z$-scores were tested and rejected on evidence.} The real corpus has
$\max|z| \in [0.2, 0.9]$, while pure Poisson noise produces $|z|$ with a median
of 1.4--1.7 and a 95th percentile of 3.0--4.0 --- it scores noise higher than
signal. The Poisson variance assumption fails because one author posting eight
times, or one news item being reshared, breaks independence. Permutation testing
assumes no distribution and is valid where the $z$-score is not.

\textbf{How much more data is needed: about 8$\times$.} Measured, not assumed, by
replicating the author pool and re-running the test --- scaling matches
$\sqrt{k}$, and the effect becomes detectable at $\approx$1{,}264 relevant posts
per window versus 158 today. Reachable via $\sim$56-day windows ($\approx$4
months) or 8$\times$ the volume; parameter tuning cannot substitute.

% ==========================================================================
\section{Two structural findings}

\textbf{The labeler was substantially wrong, and nobody had checked.} A stratified
100-post sample with Wilson intervals put \code{is\_relevant} precision at
\textbf{56.0\%} --- nearly half of what the system called Canada-related
cybersecurity discussion was not. A corroboration rule (weak abbreviations count
only when an independent Canada-side signal also fires) lifted it to
\textbf{86.3\%}, recall 85.6\% $\to$ \textbf{91.7\%}. The pilot headline moved
from 27/day to 21.3/day. The fix was not a better rule --- it was measuring the
rule at all.

\textbf{A quarter of the corpus has no audience.} \textbf{11 of 372 accounts
(3.0\%) produce 158 of 650 relevant posts (24.3\%)}, with mean engagement
\textbf{0.20} against \textbf{5.34} for everyone else. They are automated
ransomware-leak-site monitors and security-conference promotion accounts.
Concretely: of 45 \code{ransomware} posts in 14 days, \textbf{40 came from 5
automated feeds}; excluding them removes the topic from the ranking entirely.
Every \code{ransomware} rise or fall in earlier reports was measuring feed
uptime. Detection is behavioral rather than a handle blocklist, and controls for
collection lag --- \code{like\_count} is a snapshot, so without a settling-time
guard every freshly collected post would look like a bot.

% ==========================================================================
\section{Limitations}

\begin{itemize}[leftmargin=1.2em,itemsep=0.1em,topsep=0.2em]
\item \textbf{Single source.} Structural, and the largest limitation.
\item \textbf{Volume is $\sim$8$\times$ short} of supporting positive trend claims.
  Defensible descriptions and defensible negative results are what it currently
  produces.
\item \textbf{The evaluation set is not independent} --- annotated by the same
  agent that wrote the labeler, which likely inflates both precision and recall.
\item \code{other} \textbf{absorbs 44\%} of relevant posts; the topic lexicon
  does not cover half the corpus.
\item \textbf{Weekend volume is $\sim$45\% of weekday volume}, which constrains
  window choice.
\item \textbf{Engagement metrics are snapshots}, not final values; cross-mode
  comparison is invalid without a correction that does not yet exist.
\end{itemize}

% ==========================================================================
\section{Next, in priority order}

\begin{enumerate}[leftmargin=1.4em,itemsep=0.1em,topsep=0.2em]
\item \textbf{Independent annotation of the evaluation set} --- one TSV file, and
  it currently gates the credibility of every precision number above.
\item \textbf{Add Reddit} --- attacks single-source bias and the 8$\times$
  shortfall at once, and is discussion-shaped rather than broadcast-shaped.
\item Extend evaluation to incremental data; move to 28--56 day windows as the
  corpus allows; correct engagement for collection lag.
\item \textbf{Only then} entities, keywords and semantic drift --- once there is
  enough data for them to mean anything.
\end{enumerate}

\vspace{0.3em}
\begin{keybox}
The most useful thing this project produced is \textbf{the ability to tell that
it has not found a trend.} A version reporting \code{data\_breach} $+56\%$ would
look more successful and would be wrong. The null model, power analysis,
broadcast filter and evaluation set are the actual deliverable --- the corpus
will grow, but the discipline to know when it is big enough had to come first.
\end{keybox}

\end{document}
